% !TEX TS-program = lualatex

\newcommand{\seriesTomeEditionNo}{1}
\newcommand{\seriesTomeNo}{1}

% For including core web pages (trim: left, bottom, right, top)
\newcommand{\bookWeb}[3]{\newpage\noindent%
\includegraphics[trim=18.16mm 20mm 10mm 18.3mm, page=#3, scale=1.065]%
{\txex s/#1/#1#2.pdf}}


% Title as text (not used so far)
\newcommand{\seriesTomeTitleText}{\en{Space Time Elements}\de{Raum Zeit Elemente}}

% Title in color as used on title pages
\newcommand{\seriesTomeTitle}{%
\begin{center}{\huge
\en{\seriesWhite{Space}\hspace{-0.02em}
\seriesRed{Time}\hspace{0.02em}
\seriesBlack{Elements}}%
\de{\seriesWhite{Raum}\hspace{0.03em}
\seriesRed{Zeit}\hspace{0.04em}
\seriesBlack{Elemente}}}
\end{center}}

% TODO ISBN correct?
%\newcommand{\seriesTomeEditionISBN}{\en{978-3-906914-06-0}\de{978-3-906914-08-4}}
\newcommand{\seriesTomeEditionISBN}{{\color{red}ONLINE SNAPSHOT \DTMnow}}

% TODO year(s) correct?
% TODO definition correct?
% Year or year span (20xx-20yy) of copyright for tome edtition,
% i.e. from first edition to current edition
%\newcommand{\seriesTomeEditionYearsCopyright}{2027}
\newcommand{\seriesTomeEditionYearsCopyright}{2025-\en{now}\de{jetzt}}

% TODO year correct?
% TODO definition correct?
% Year the current edition (identified via ISBN) was published
%\newcommand{\seriesTomeEditionYearFirstPublished}{2027}
\newcommand{\seriesTomeEditionYearFirstPublished}{\the\year}

\seriesFrontCover

\seriesFirstTitlePage
\seriesFirstTitlePageReverse
\seriesSecondTitlePage
\seriesSecondTitlePageReverse
\seriesTableOfContents

%\seriesEmptyPage
\seriesAboutSection

%••••••••••••••••••••••••••••••••••••••••••••••••••••••••••••••••••••••••••••••••••••••••••••••••••••
\section{\en{Preface}\de{Vorwort}}
%••••••••••••••••••••••••••••••••••••••••••••••••••••••••••••••••••••••••••••••••••••••••••••••••••••

\en{…}%
\de{…}

%\en{My first contact with philosophy was when I was around ten years old.}%
%\de{Mein erster Kontakt mit Philosophie war als ich etwa zehn Jahre alt war.}
%%
%\en{Me and my father were riding on a chairlift in winter
%and my father told me Plato's \textsl{Allegory of the Cave}
%with ants sitting in their earth
%and just seeing the shadows of the chairs of the chairlift
%cast into the snow in front of them.}%
%\de{Ich und mein Vater fuhren auf einem Sessellift im Winter
%und mein Vater erzählte mir Platons \textsl{Höhlengleichnis}
%mit Ameisen, die in ihrem Bau sind
%und nur die Schatten der Sessel des Sessellifts sehen,
%die vor ihnen in den Schnee geworfen werden.}
%%
%\en{What would they imagine the world to be
%with shadows of skiers on chairs rolling by,
%one after the other?}%
%\de{Wie würden sie sich die Welt vorstellen,
%mit Schatten von Skifahrern auf Sesseln,
%die einer nach dem anderen vorbeiziehen würden?}
%
%\en{I guess that relativizes almost anything you could write in any book…}%
%\de{Ich würde sagen, das relativiert fast alles, was man je in irgendeinem Buch schreiben könnte…}
%
%\en{Still, there are very interesting new ideas in this book,
%mostly pragmatic new approaches to questions
%that were considered since at least the ancient Greeks.}%
%\de{Und doch hat es sehr interessante neue Ideen in diesem Buch,
%vorwiegend pragmatische neue Zugänge zu Fragen,
%die schon seit mindestens den alten Griechen betrachtet wurden.}
%
%\en{The immediate perception of space and time in personal experience
%seems to lead to something that is very similar to
%especially Aristotle’s views on the classical elements
%Fire, Air, Water and Earth, plus the fifth, Ether.}%
%\de{Die unmittelbare Wahrnehmung von Raum und Zeit in der persönlichen Erfahrung
%scheint zu etwas zu führen,
%das insbesondere Aristoteles’ Ansichten über die klassischen Elemente
%Feuer, Luft, Wasser und Erde sowie dem fünften, dem Äther,
%sehr ähnlich ist.}
%
%\en{And things would seem to even go way beyond that.}%
%\de{Und es schiene sogar noch viel weiter zu gehen.}
%
%\en{But so far this is largely just what I found over the years.}%
%\de{Aber soweit ist das weitgehend nur, was ich über die Jahre gefunden habe.}
%%
%\en{Wikipedia would categorize it as “original work”,
%while I am still a physicist with a PhD,
%so personally I do see it as relevant research in many ways.}%
%\de{Wikipedia würde es als “original work” kategorisieren,
%wohingegen ich immer noch promovierter Physiker hin,
%also persönlich es als in vielerlei Hinsicht relevante Forschung erachte.}
%
%\en{But see for yourselves on the following pages.}%
%\de{Aber seht selbst auf den folgenden Seiten.}
%%
%\en{And remember:
%This is just a first glimpse at something very interesting
%with maybe lots of potential,
%but not yet a closed and verified scientific theory.}%
%\de{Und vergesst nicht:
%Soweit ist das hier erst ein erster Einblick in etwas sehr Interessantes
%mit vielleicht sehr viel Potential,
%aber noch keine abgeschlossene und verifizierte wissenschaftliche Theorie.}
%
%\vspace{3mm}
%\noindent
%\en{Adliswil (near Zürich), 1 June 2025\newline}%
%\de{Adliswil (nahe Zürich), 1.~Juni 2025\newline}
%Alain Stalder


%••••••••••••••••••••••••••••••••••••••••••••••••••••••••••••••••••••••••••••••••••••••••••••••••••••
%\chapter{Introduction}
%••••••••••••••••••••••••••••••••••••••••••••••••••••••••••••••••••••••••••••••••••••••••••••••••••••

%\noindent
%In late 2017, Google’s DeepMind AlphaZero AI program
%beat the then strongest chess program Stockfish easily, at
%a time when humans did not have any chance any more
%anyways. But the way AlphaZero won was what was
%really astonishing. Unlike chess programs before it, it had
%not learned from previous games played by humans (or
%machines), it had learned strictly only from the chess rules.
%
%This book has quite a similar approach to reality and the
%world, only so far not with AI, but with good old human
%reasoning. The starting point here is a completely fresh
%look at the world, one that does not assume anything about
%the world to be known, as if just having opened one’s eyes
%(and all other senses) to the world for the first time ever
%and trying to make sense of all that is going on in what
%seems to be an ‘outside‘ and an ‘inside’.
%
%AlphaZero made moves nobody expected, sometimes
%even sacrificing the strongest piece on the board, the queen,
%which humans almost never do. Similarly, even though
%the findings here usually closely mirror previous findings of
%humans about life and the world that have been acquired
%across millenia, sometimes they differ, as follows…
%
%Since the approach here is strictly scientific in the sense
%that it all starts with observation and tries to be as logically
%consistent as ever possible, it is never in contradiction with
%any verifiable facts that science can regularly predict.
%
%But the approach also naturally leads to concepts that
%are not directly part of contemporary science. For example,
%one of the first concepts that appears is that there would be
%something that comes in 4+1 ‘parts’ and has astonishing
%similarities with ancient ideas about 4+1 elements (“fire,
%air, water, earth” and “ether”), or also the similar concept
%of eight trigrams of the Chinese I Ching, only in both cases
%in a more modern and well-defined way. And also emerging
%quickly is a more modern view of something that was often
%called and perceived as gods and godesses in the past, or
%similar concepts like a collective unconscious or a world
%soul, only here explained more scientifically as unconscious
%collective connections and effects.
%
%It is important to note that the two things I just mentioned,
%concepts that resemble ancient elements and deities,
%seem to follow quite naturally and quite quickly from that
%fresh look at the world, about as astonishingly as some of
%AlphaZero’s moves. Of course, so far (in the first part of
%the 21st century), it has just been me doing this, so my
%approach might be biased or overlook other possibilities,
%but it clearly shows that there is more to be found with
%that approach. So, if you are looking for something novel
%to do, this might be a really good opportunity\,!

%••••••••••••••••••••••••••••••••••••••••••••••••••••••••••••••••••••••••••••••••••••••••••••••••••••
%\en{\chapter[Elements as Space and Time]{Elements as\newline Space and Time}}%
%\de{\chapter[Elemente als Raum und Zeit]{Elemente als\newline Raum und Zeit}}
%••••••••••••••••••••••••••••••••••••••••••••••••••••••••••••••••••••••••••••••••••••••••••••••••••••

%••••••••••••••••••••••••••••••••••••••••••••••••••••••••••••••••••••••••••••••••••••••••••••••••••••
%\en{\chapter[Deities as Collective Effects]{Deities as\newline Collective Effects}}%
%\de{\chapter[Gottheiten als kollektive Effekte]{Gottheiten als\newline kollektive Effekte}}
%••••••••••••••••••••••••••••••••••••••••••••••••••••••••••••••••••••••••••••••••••••••••••••••••••••

%••••••••••••••••••••••••••••••••••••••••••••••••••••••••••••••••••••••••••••••••••••••••••••••••••••
%\en{\chapter[Astrology Void of Stars]{Astrology\newline Void of Stars}}%
%\de{\chapter[Astrologie ohne Sterne]{Astrologie\newline ohne Sterne}}
%••••••••••••••••••••••••••••••••••••••••••••••••••••••••••••••••••••••••••••••••••••••••••••••••••••

%••••••••••••••••••••••••••••••••••••••••••••••••••••••••••••••••••••••••••••••••••••••••••••••••••••
%\en{\chapter[Miscellaneous ‘Phenomena’]{Miscellaneous\newline ‘Phenomena’}}%
%\de{\chapter[Verschiedene ‘Phänomene’]{Verschiedene\newline ‘Phänomene’}}
%••••••••••••••••••••••••••••••••••••••••••••••••••••••••••••••••••••••••••••••••••••••••••••••••••••

%••••••••••••••••••••••••••••••••••••••••••••••••••••••••••••••••••••••••••••••••••••••••••••••••••••
\addcontentsline{toc}{chapter}{\en{Appendix}\de{Anhang}: exactphilosophy.net}
\chapter*{\en{Appendix}\de{Anhang}:\newline{exactphilosophy.net}}%
%••••••••••••••••••••••••••••••••••••••••••••••••••••••••••••••••••••••••••••••••••••••••••••••••••••

\en{This chapter reproduces the core content of my website
\href{https://www.exactphilosophy.net/}{\color{seriesXphi}{exactphilosophy.net}}
at the time of writing.}%
\de{Dieses Kapitel gibt den Kerninhalt meiner Website
\href{https://www.exactphilosophy.net/}{\color{seriesXphi}{exactphilosophy.net}}
zur Zeit des Schreibens wieder.}

\en{The style is very minimalist,
both from a scientific and an artistic perspective,
in the hope to come close to nature.}%
\de{Der Stil ist sehr minimalistisch,
sowohl aus wissenschaftlicher wie auch aus künstlerischer Perspektive,
in der Hoffnung,
der Natur nahe zu kommen.}

\en{It is the fruit of countless hours since at least 2005.}%
\de{Der englische Text ist die Frucht ungezählter Stunden seit mindestens 2005.
Die deutsche Übersetzung,
die im Wesentlichen im Juli 2023 innerhalb weniger Tage entstand,
ist auch schön,
wird aber wohl nie ganz so schön sein können,
auch wenn sie ganz eigene schöne Eigenschaften hat.}

%••••••••••••••••••••••••••••••••••••••••••••••••••••••••••••••••••••••••••••••••••••••••••••••••••••
% Include web pages
%••••••••••••••••••••••••••••••••••••••••••••••••••••••••••••••••••••••••••••••••••••••••••••••••••••

\newcommand{\toctypewelcome}{section}
\newcommand{\toctitlewelcome}{\en{Welcome}\de{Willkommen}}

\newcommand{\toctypeway}{section}
\newcommand{\toctitleway}{\en{way}\de{weg}}
\newcommand{\toctypespacetime}{subsection}
\newcommand{\toctitlespacetime}{\en{space and time}\de{raum und zeit}}
\newcommand{\toctypemetamorphosis}{subsection}
\newcommand{\toctitlemetamorphosis}{\en{metamorphosis}\de{metamorphose}}
\newcommand{\toctypegreek}{subsection}
\newcommand{\toctitlegreek}{\en{greek philosophy}\de{griechische philosophie}}
\newcommand{\toctypeiching}{subsection}
\newcommand{\toctitleiching}{\en{i ching}\de{i ging}}
\newcommand{\toctypeorigins}{subsection}
\newcommand{\toctitleorigins}{\en{origins}\de{ursprünge}}

\newcommand{\toctypeevolutions}{section}
\newcommand{\toctitleevolutions}{\en{evolutions}\de{evolutionen}}
\newcommand{\toctypefeelings}{subsection}
\newcommand{\toctitlefeelings}{\en{mixed feelings}\de{gemischte gefühle}}
\newcommand{\toctypestarsigns}{subsection}
\newcommand{\toctitlestarsigns}{\en{star signs}\de{sternzeichen}}
\newcommand{\toctypeartemis}{subsection}
\newcommand{\toctitleartemis}{\en{artemis}\de{artemis}}

\newcommand{\toctypesynthesis}{section}
\newcommand{\toctitlesynthesis}{\en{synthesis}\de{synthese}}
\newcommand{\toctypemeasurement}{subsection}
\newcommand{\toctitlemeasurement}{\en{measurement}\de{messung}}
\newcommand{\toctypepsyche}{subsection}
\newcommand{\toctitlepsyche}{\en{psyche}\de{psyche}}
\newcommand{\toctypedreams}{subsection}
\newcommand{\toctitledreams}{\en{dreams}\de{träume}}

\newcommand{\toctypereferences}{section}
\newcommand{\toctitlereferences}{\en{References}\de{Referenzen}}

\newcommand{\toctypeseries}{section}
\newcommand{\toctitleseries}{\en{Series}\de{Reihe}}

\newcommand{\toctypelinks}{section}
\newcommand{\toctitlelinks}{\en{Links}\de{Links}}

\en{\input{\txex s/web-en.inc}}%
\de{\input{\txex s/web-de.inc}}

\seriesFinalPageArtecatLogo
% no empty pages at the end needed; ingramspark adds
% 2 or 3 pages to make even plus printer's page with bar code;
% just add also 2 or 3 pages when getting the cover template.

\seriesBackCover
